\section{A complete local obstruction for the identity cubic unit}
\label{cl-section}

The identity unit \(u=1\) must be distinguished from \(u=\zeta\).
The explicit cubic quotient in Section~\ref{hg-section} gives a
complete local exclusion for the former, while leaving the two
nonidentity unit classes open.

\begin{theorem}[The projective cubic quotient has no \(3\)-adic point]
\label{cl-no-local-point}
Put
\begin{align*}
 C(S,T)&=S^3+3S^2T-T^3,\\
 E_3(S,T)&=S^3-S^2T-4ST^2-T^3,\qquad D(S,T)=-3E_3(S,T).
\end{align*}
Let \(H\) be the smooth projective double cover
\begin{equation}\label{cl-sextic}
 Y^2=C(S,T)D(S,T),
\end{equation}
where \(Y\) has weight \(3\). Then \(H(\mathbb Q_3)=\varnothing\).
This is the third quotient in Theorem~\ref{hg-three-quotients}
at \(g=3\), \(u=1\).
\end{theorem}
\begin{proof}
The binary sextic has six simple branch points by the proved HG
calculation. A rational point maps to \(\mathbb P^1(\mathbb Q_3)\).
Choose \(S,T\in\mathbb Z_3\) with \(\min(v_3(S),v_3(T))=0\).
Its fibre would supply a \(Y\in\mathbb Q_3\) satisfying
\[
 Y^2=-3C(S,T)E_3(S,T).
\]
This includes infinity, represented by \((1,0)\). Scaling the base
coordinates changes the right side by a sixth power, so the square
criterion does not depend on that choice.

On the four residue points represented by
\((1,0),(0,1),(1,1),(-1,1)\), the values of \(E_3\) are respectively
\(1,-1,1,1\) modulo \(3\). Thus \(E_3(S,T)\) is always a unit.
If \(S\not\equiv T\pmod3\), then
\(C(S,T)\equiv(S-T)^3\) is a unit as well. The right side has
valuation \(1\), so cannot be a square.

Otherwise \(T\) is a unit and \(S=T+3K\), with \(K\in\mathbb Z_3\).
The exact expansion is
\begin{equation}\label{cl-shift}
 C(T+3K,T)=3(T^3+9KT^2+18K^2T+9K^3).
\end{equation}
Hence \(v_3(C)=1\), \(C/3\equiv T^3\pmod3\), and also
\(E_3\equiv T^3\pmod3\). The right side has valuation \(2\), but
after division by \(9\) its unit part is
\(-T^6\equiv-1\pmod3\). A \(3\)-adic unit square has residue \(1\).
This is another contradiction, covering every projective fibre.
\end{proof}

\begin{corollary}[Propagation to identity-unit exponents divisible by three]
\label{cl-propagation}
For every odd positive integer \(g\) divisible by \(3\), the curves
of Section~\ref{rl-section} satisfy
\[
 D_{g,1}(\mathbb Q_3)=C_{g,1}(\mathbb Q_3)=\varnothing.
\]
Their entire \(\mathbb Q\)-rational loci are therefore empty.
\end{corollary}
\begin{proof}
For \(g=3\), the quotient morphism
\(D_{3,1}\to H\) from Theorem~\ref{hg-three-quotients} and
Theorem~\ref{cl-no-local-point} give the claim.
For \(g=3k\), successive powers give the identity of projective maps
\[
 L_{1,g}=L_{1,3}\circ L_{1,k}.
\]
On the generic fibre, \((t,s)\mapsto(t,L_{1,k}(s))\) defines a
nonconstant rational map \(D_{g,1}\dashrightarrow D_{3,1}\).
The curve extension theorem~\cite{StacksCurves2026} extends it
over the smooth projective source, including all infinite and
zero-denominator chart points. It is defined over \(\mathbb Q\),
hence over \(\mathbb Q_3\). A local point would map to a forbidden
point of \(D_{3,1}\). The forgetful morphism
\(C_{g,1}\to D_{g,1}\) proves the assertion for the higher cover.
\end{proof}

For an actual primitive seed with \(M=a^2+ab+b^2=U^2\),
the elementary norm calculation already gives \(3\nmid M\):
its primitive \(3\)-adic depth is at most one and must be even.
An identity-unit integral cube \((S+T\zeta)^3\), however, has
\(\zeta\)-coefficient \(3ST(S+T)\), so cannot equal \(ab+M\zeta\).
Thus this integral subcase already has an elementary obstruction.
The new curve statement records the full local obstruction,
including nonpositive and infinite charts, and its projective
propagation. It does not exclude \(u=\zeta,\zeta^2\), moving
residuals, or exponents not divisible by three.

The six declarations of
\texttt{CubicUnitLocal\allowbreak Arithmetic.lean} prove the
actual shift identity, modular reduction of the sextic, the full
kernel-checked \(\operatorname{Fin}(27)\) table, its transfer to
all integer coordinates, and the impossibility of an integer
square at primitive integer coordinates. The finite table has
no axiom dependencies; the other queries use only the standard
Lean axioms. These arithmetic statements do not formalize
\(\mathbb Q_3\), the projective curves, or the propagation morphism.
